\chapter{Introdução}
\label{cap:introducao}

\section{Motivação e Justificativa}

Esta seção deve mostrar que o trabalho aqui descrito tem valor acadêmico e prático. 
Apresente inicialmente o cenário geral do trabalho.
Explique brevemente os conceitos do trabalho realizado, que é geralmente relacionado ao desenvolvimento ou manutenção de software, e por que esses processos são relevantes para a indústria de software e para a sociedade em geral.

Explique o problema que motivou o desenvolvimento ou a manutenção do software em que você atuou.
Mostre a importância de resolver este tipo de problema com atuação em equipe, destacando desafios comuns e a necessidade de boas práticas de processo para organizar o trabalho da equipe.
Justifique por que relatar essa experiência traz contribuições acadêmicas e práticas. Argumente, por exemplo:

\begin{itemize}
    \item dependência crescente da sociedade atual em relação a software.
    \item relevância de processos eficientes de desenvolvimento e manutenção para a indústria de software.
    \item importância da experiência prática para formação profissional.
    \item contribuição do relato para outros estudantes, equipes ou organizações.
\end{itemize}

Mostre por que a experiência relatada neste trabalho é relevante para estudo e prática de Engenharia de Software, por exemplo, em termos de desafios enfrentados pela equipe, complexidade do sistema ou importância do produto para as partes interessadas (usuários, organização ou toda a sociedade).

\section{Contexto da Experiência}

Situe o leitor no tema do seu projeto específico, tanto em termos de objetivo do projeto quanto de organização do processo de trabalho e das equipes. Descreva: 

\begin{itemize}
    \item o ambiente onde a experiência ocorreu: empresa, laboratório, projeto acadêmico ou organização.
    \item tipo de equipe (tamanho, papéis, metodologia, subdivisões).
    \item contexto do projeto (novo desenvolvimento, manutenção evolutiva, correção de bugs etc.).
\end{itemize}

Não entre em detalhes técnicos. Apenas apresente o cenário geral do projeto. Detalhes aparecem no Capítulo~\ref{cap:relato}.

\section{Objetivo do Trabalho}

Um trabalho acadêmico de conclusão de curso deve ter um objetivo geral, que pode ser desdobrado em objetivos específicos. Defina inicialmente um único objetivo geral para o seu TCC e desdobre este objetivo geral em objetivos menores e mais específicos, necessários para alcançar o objetivo geral.

\subsection{Objetivo Geral}

Declare claramente o objetivo principal do TCC.
Por exemplo: "O objetivo deste Trabalho de Conclusão de Curso (TCC) é descrever e analisar a experiência de participação em uma equipe de desenvolvimento e manutenção de software, analisando práticas, desafios e aprendizados vivenciados nesta experiência."

\subsection{Objetivos Específicos}

Apresente os objetivos específicos que detalham o objetivo geral e representam marcos importantes para cumprir este objetivo.
Exemplos:

\begin{enumerate}
    \item Descrever o processo de trabalho adotado pela equipe.
    \item Relatar as atividades realizadas no desenvolvimento e manutenção de software.
    \item Identificar desafios técnicos e organizacionais enfrentados.
    \item Apresentar as soluções adotadas e seus resultados.
    \item Refletir sobre as lições aprendidas ao longo da experiência.
\end{enumerate}

É importante retomar estes objetivos nas conclusões (Capítulo~\ref{cap:conclusoes}) sintetizando como cada objetivo foi alcançado e referenciando as partes (seções do TCC) que apresentam a realização destes objetivos. 

\section{Metodologia do Trabalho}

Explique como o TCC foi realizado em relação ao tipo de metodologia acadêmica. No caso do curso de Engenharia de Software, a metodologia do TCC envolve tipicamente:

\begin{itemize}
    \item Estudo de caso: análise crítica feita pelo autor do TCC sobre o trabalho de Engenharia de Software feito por terceiros; ou
    \item Pesquisa-ação: resolução de problemas de Engenharia de Software feita pelo próprio autor do TCC.
\end{itemize}

A maioria dos TCCs seguem a metodologia de pesquisa-ação, mas alguns TCCs envolvem também um estudo de caso que identifica problemas a serem tratados com a metodologia pesquisa-ação. Apresente aqui uma visão geral dos conceitos da metodologia acadêmica utilizada no TCC.

Não confunda a \textbf{metodologia de desenvolvimento do trabalho acadêmico} (pesquisa-ação, estudo de caso) com a \textbf{metodologia de desenvolvimento ou manutenção de software} usada pela equipe de Engenharia de Software (Scrum, Kanban etc.).
Esta metodologia de Engenharia de Software será detalhada no Capítulo~\ref{cap:relato}.

Descreva as fontes bibliográficas utilizadas na realização do trabalho, inclusive as fontes que explicam a metodologia acadêmica aplicada. Exemplos de fontes bibliográficas: literatura técnico-científica, documentação do projeto, sistemas similares, documentação de ferramentas, etc.

\section{Organização do Trabalho}

Apresente um resumo do conteúdo de cada capítulo subsequente para ajudar o leitor a entender o fluxo do documento e onde se localiza cada tópico do trabalho.
Exemplo:

\begin{itemize}
    \item Capítulo 2: Bases teóricas e tecnológicas utilizadas nas atividades realizadas.
    
    \item Capítulo 3: Descrição da experiência individual vivenciada e do trabalho em equipe, com discussão dos resultados e lições aprendidas.
    
    \item Capítulo 5: Considerações finais, comparação com trabalhos correlatos e indicações de trabalhos futuros.
    
\end{itemize}
